%--- iliad ---
% separateSolutions: true
% title: Condensation — solutions
% summary: >-
%   Complete solutions to exercises E1–E10 and the correspondence proof.
% unlisted: true
%--- end ---
\documentclass[11pt]{article}
\IfFileExists{iliad.sty}{\usepackage{iliad}}{\usepackage{../iliad}}
\geometry{a4paper,margin=25mm}
\usepackage{microtype}
\newcommand{\PI}{\mathcal{P}^{+}(I)}
\newcommand{\Up}{\mathord{\supseteq}}
\title{Condensation: Solutions}
\author{Satya Benson}
\date{21 September 2026}
\begin{document}
\maketitle

\section{Prerequisites}
This is the separate answer key for \emph{Condensation}. Exercise identifiers
E1--E10 match the student worksheet. All variables have finite ranges, all logs
have base two, and unused latent slots are constant. Write
$\Up A=\{C:A\subseteq C\}$,
and $\mathcal{J}_B=\{C:C\cap B\ne\varnothing\}$.
A tuple indexed by a family $\mathcal{F}$ is written $Y_{\mathcal{F}}$.

\begin{learningoutcomes}
\begin{itemize}
\item Check all example calculations and identify exactly where a condition fails.
\item Follow a complete proof of the exact and quantitative correspondence result.
\item Extend the argument to observation blocks and verify the conditioned-score identity.
\end{itemize}
\end{learningoutcomes}

\section{E1: Overlapping bits}
\label{sec:e1}
The independent fair bits are $S,P,Q$, and
$X_1=(S,P)$, $X_2=(S,Q)$, $X_3=(P,Q)$.
\begin{enumerate}[label=(\alph*)]
\item $H(X_1)=2$. The tuple of all observations determines $(S,P,Q)$ and is
determined by it, so $H(X_1,X_2,X_3)=3$.
The pair $(X_1,X_2)$ already determines all three bits, hence
\[
 I(X_1;X_2)=H(X_1)+H(X_2)-H(X_1,X_2)=2+2-3=1.
\]
Its shared bit is $S$.
\item The families, including constant slots, are
\[
\begin{split}
 \Up\{1\}&=\{\{1\},\{1,2\},\{1,3\},\{1,2,3\}\},\\
 \Up\{2\}&=\{\{2\},\{1,2\},\{2,3\},\{1,2,3\}\},\\
 \Up\{1,2\}&=\{\{1,2\},\{1,2,3\}\}.
\end{split}
\]
The first tuple has nonconstant entries $S,P$, so determines $X_1$.
Intersecting the first two families leaves $\{\{1,2\},\{1,2,3\}\}=\Up\{1,2\}$,
the third family listed, whose only nonconstant content is $S$.
\item Recover $X_1$ by taking the ordered pair
$(Y_{\{1,2\}},Y_{\{1,3\}})$. Recover $Y_{\{1,2\}}$ from $X_1$ by
projecting onto its first coordinate.
\item Addresses $\{1,2\}$ and $\{1,3\}$ overlap but neither contains the
other. Thus the contribution pattern need not form a nested tree.
\end{enumerate}
The original model also satisfies the Markov condition. All nonconstant slots
are mutually independent, so disjoint subcollections are independent, including
after conditioning on the remaining common subcollection. Reconstruction follows from the indicated coordinate projections.

\section{E2: Two failed models}
\label{sec:e2}
\begin{enumerate}[label=(\alph*)]
\item In the singleton-copy model $Y_{\{i\}}=X_i$, the contributing family
for $i$ includes $X_i$ itself. Thus LVM holds. Each
nonconstant latent has just one contributing observation, from which it is
recovered identically, so reconstruction holds too.

For $\mathcal{F}=\Up\{1\}$ and $\mathcal{G}=\Up\{2\}$, all entries
of $Y_{\mathcal{F}\cap\mathcal{G}}$ are constant. The only nonconstant entries
of the two larger tuples are $X_1$ and $X_2$. Consequently
\[
 I(Y_{\mathcal{F}};Y_{\mathcal{G}}\mid Y_{\mathcal{F}\cap\mathcal{G}})
 =I(X_1;X_2)=1.
\]
Both singleton slots determine $S$, while their common ancestor slots are
constant. This violates the Markov condition. It is enough to
exhibit this one pair of upward-closed families; finding the maximum defect
is unnecessary.
\item With only $Y_{\{1,2,3\}}=(S,P,Q)$ nonconstant, each observation is
obtained by projection, so LVM holds. Every nonempty
upward-closed family contains the top address. If both families are nonempty,
their intersection contains that top slot and determines all latent variables;
the conditional mutual information is zero. If either is empty it is zero as
well. Thus the Markov condition holds.

But $X_1=(S,P)$ leaves the independent fair bit $Q$ unknown, so
\[
 H(Y_{\{1,2,3\}}\mid X_1)=H(Q\mid S,P)=1.
\]
Reconstruction fails.
\item LVM requires enough latent information to determine
each observation. Reconstruction excludes information placed at a shared
address that an individual contributing observation cannot recover. The Markov
condition excludes dependence between address collections not accounted for
by their common slots. The singleton model satisfies the second requirement but
fails the third; the top-slot model satisfies the third but fails the second.
\end{enumerate}

\subsection{Check on the existence example}
The existence subsection in the student sheet is a proved example, not an
additional exercise. For a binary pair with full support, a top slot determined
by each observable must be constant: writing it as $f(X_1)=g(X_2)$, each of the
four pairs imposes an equality between the corresponding function values.
With a constant top, LVM makes $X_1$ a function of
$Y_{\{1\}}$ and $X_2$ a function of $Y_{\{2\}}$. Thus data processing yields
\[
 I(X_1;X_2)\le I(Y_{\{1\}};Y_{\{2\}})=\eta_Y.
\]
The equality holds because, on the two-observable address poset, the only
pair of incomparable nonempty upward-closed families consists of the two
singleton towers. Every other Markov term is zero. For the noisy fair-bit
pair, $H(X_2)=1$ and $H(X_2\mid X_1)=h_2(q)$, proving the stated lower bound.
It is attained by the singleton-copy model $Y_{\{i\}}=X_i$, with constant top.
This model has both exact reconstruction directions but a positive Markov
defect. Setting that defect to zero is impossible. The argument does not bound
the tradeoff when reconstruction is allowed to be inexact.

\section{E3: Collections, not individual slots}
\label{sec:e3}
Here $C,R,T$ are independent fair bits. In model $Y$ the three slots are
$C,R,T$; in model $Z$ they are $C,(C,R),(C,T)$.
\begin{enumerate}[label=(\alph*)]
\item Both contributing families determine their observations. All slots are
recoverable from each observation to which they contribute. For example,
$C$ is the first coordinate of either $X_1$ or $X_2$, and $Z_{\{1\}}=X_1$.
For two observables the only nontrivial Markov requirement is independence of
the two singleton slots given the top slot. For $Y$ this is
$I(R;T\mid C)=0$. For $Z$ it is
\[
 I((C,R);(C,T)\mid C)=I(R;T\mid C)=0.
\]
Both are therefore perfect condensations.
\item $H(Z_{\{1\}}\mid Y_{\{1\}})=H(C,R\mid R)=H(C)=1$.
\item The $Y$ tower at $\{1\}$ consists of $(R,C)$, and the $Z$ tower
consists of $((C,R),C)$. Each is a deterministic function of the other. A copy
of the higher-address information can appear in a lower slot without changing
the tower's total information. Thus the individual slots need not correspond
even when the relevant towers do.
\end{enumerate}

\section{E4: The two-witness lemma}
\label{sec:e4}
For arbitrary finite $U,M,P,Q$, the exact identity is
\begin{align*}
 &H(U\mid M,P)+H(U\mid M,Q)+I(P;Q\mid M)-H(U\mid M)\\
 &\hspace{15mm}=I(P;Q\mid U,M)+H(U\mid M,P,Q).
\end{align*}
For a direct verification, first subtract the two conditional mutual
informations using their entropy definitions:
\begin{align*}
 I(P;Q\mid M)-I(P;Q\mid U,M)
 &=H(P\mid M)+H(Q\mid M)-H(P,Q\mid M)\\
 &\quad-H(P\mid U,M)-H(Q\mid U,M)+H(P,Q\mid U,M)\\
 &=I(U;P\mid M)+I(U;Q\mid M)-I(U;P,Q\mid M)\\
 &=H(U\mid M)-H(U\mid P,M)-H(U\mid Q,M)\\
 &\quad+H(U\mid P,Q,M).
\end{align*}
Rearrangement gives the displayed slack identity. Its right side is
nonnegative, proving the lemma.

If $H(U\mid M,P)=H(U\mid M,Q)=I(P;Q\mid M)=0$, the inequality and
nonnegativity imply $H(U\mid M)=0$. Thus a target recoverable from each
witness is recoverable from their designated common part when the witnesses
are conditionally independent given that part.

The first two terms measure what each witness, together with $M$, still fails
to determine about $U$. The third measures dependence between $P,Q$ not
accounted for by $M$. No independence assumption is needed for the inequality
itself; exact independence is the case where the third error is zero.

\section{E5: The dependence term is needed}
\label{sec:e5}
With $U=P=Q$ a fair bit and $M$ constant,
\[
 H(U\mid M)=1,\quad H(U\mid M,P)=H(U\mid M,Q)=0,
 \quad I(P;Q\mid M)=1.
\]
The inequality is the equality $1\le0+0+1$. Without its mutual-information
term it would assert $1\le0$. The two witnesses each determine the same bit while
the designated common variable is constant, as in the singleton-copy model
of E2.

\section{E6: Transfer through an observable}
\label{sec:e6}
If $X_i$ is a function of $Z_{\Up\{i\}}$ almost surely, adding $X_i$ to
that conditioning tuple changes no information. Hence
\[
 H(U\mid Z_{\Up\{i\}})
 =H(U\mid Z_{\Up\{i\}},X_i)
 \le H(U\mid X_i).
\]
The inequality is conditioning monotonicity. It does not assume that $U$ and
the receiver tuple are conditionally independent given $X_i$. They can share
additional information beyond $X_i$.

\section{E7: Intersections of address families}
\label{sec:e7}
\begin{enumerate}[label=(\alph*)]
\item The first two families are listed in E1, and
\[
 \Up\{3\}=\{\{3\},\{1,3\},\{2,3\},\{1,2,3\}\}.
\]
Thus $\Up\{1\}\cap\Up\{2\}=\{\{1,2\},\{1,2,3\}\}$,
and intersecting with $\Up\{3\}$ leaves only $\{1,2,3\}$.
\item If $C$ belongs to every family in an intersection and $D\supseteq C$,
upward closure places $D$ in each such family, hence in the intersection.
\item For an address $C$,
\[
 C\in\bigcap_{i\in A}\Up\{i\}
 \quad\Longleftrightarrow\quad
 (\forall i\in A)\ i\in C
 \quad\Longleftrightarrow\quad A\subseteq C
 \quad\Longleftrightarrow\quad C\in\Up A.
\]
\item Put $M=Z_{\mathcal{F}\cap\mathcal{G}}$,
$P=Z_{\mathcal{F}}$, $Q=Z_{\mathcal{G}}$. The tuple $M$ is a subtuple of
both $P$ and $Q$. Thus $H(U\mid M,P)=H(U\mid P)$ and similarly for $Q$.
By the definition of $\eta_Z$, the lemma gives
\[
 H(U\mid Z_{\mathcal{F}\cap\mathcal{G}})
 \le H(U\mid Z_{\mathcal{F}})+H(U\mid Z_{\mathcal{G}})+\eta_Z.
\]
Upward closure is needed to invoke the bound $\eta_Z$.
\end{enumerate}

\section{E8: Complete correspondence proof}
\label{sec:e8}
Fix a nonempty $A=\{i_1,\ldots,i_k\}$, put
$\mathcal{G}_j=\bigcap_{\ell=1}^{j}\Up\{i_\ell\}$.
Every $\Up\{i\}$ is upward closed, as is every $\mathcal{G}_j$ by E7.
We prove
\begin{equation}
 H(Y_{\Up A}\mid Z_{\mathcal{G}_j})\le
 \sum_{\ell=1}^{j}H(Y_{\Up A}\mid X_{i_\ell})+(j-1)\eta_Z
 \label{eq:induction-sol}
\end{equation}
for $1\le j\le k$.

For $j=1$, $\mathcal{G}_1=\Up\{i_1\}$, and E6 gives
$H(Y_{\Up A}\mid Z_{\mathcal{G}_1})\le H(Y_{\Up A}\mid X_{i_1})$. This is the claimed
bound with zero intersection terms.

Suppose \eqref{eq:induction-sol} holds for $j-1$, where $2\le j\le k$.
Apply E7 to the upward-closed families $\mathcal{G}_{j-1}$ and
$\Up\{i_j\}$. Their intersection is $\mathcal{G}_j$, so
\begin{align*}
 H(Y_{\Up A}\mid Z_{\mathcal{G}_j})
 &\le H(Y_{\Up A}\mid Z_{\mathcal{G}_{j-1}})
       +H(Y_{\Up A}\mid Z_{\Up\{i_j\}})+\eta_Z\\
 &\le \sum_{\ell=1}^{j-1}H(Y_{\Up A}\mid X_{i_\ell})+(j-2)\eta_Z
       +H(Y_{\Up A}\mid X_{i_j})+\eta_Z\\
 &=\sum_{\ell=1}^{j}H(Y_{\Up A}\mid X_{i_\ell})+(j-1)\eta_Z.
\end{align*}
The second line uses the induction hypothesis and E6. This completes the
induction. At $j=k$, E7 gives $\mathcal{G}_k=\Up A$, hence
\[
 H(Y_{\Up A}\mid Z_{\Up A})
 \le\sum_{i\in A}H(Y_{\Up A}\mid X_i)+(|A|-1)\eta_Z.
\]
For $|A|=1$ this is exactly the observable-transfer inequality, so no intersection
or Markov error is required. There are $k$ reconstruction errors and $k-1$
Markov defects.
Bounding each reconstruction error by $\delta$ and each Markov defect by
$\eta$ gives
$k\delta+(k-1)\eta$, or $(2k-1)\varepsilon$ for a common budget.

If $Y,Z$ are perfect condensations, each reconstruction error is zero: for every
$C\supseteq A$ and $i\in A$, reconstruction says $Y_C$ is a function
of $X_i$. The finite tuple of such variables is consequently a function of
$X_i$. Also $\eta_Z=0$. The inequality gives zero conditional entropy in the
forward direction. Interchanging $Y,Z$, using reconstruction in $Z$
and $\eta_Y=0$, proves the reverse direction. Zero conditional entropy is
almost-sure deterministic recoverability in our finite setting.

The proof uses LVM in $Z$ for E6, the sender's reconstruction to bound those errors, and the receiver's Markov condition
at intersections. It needs no cross-model independence, and its finite
coupling preserves all marginal conditions used.

Finally, approximate slotwise errors do not directly give the same error for
a tuple. For example, if two independent fair bits $R,T$ remain unknown given
a constant observation, each conditional entropy is one bit but the tuple's
is two bits. Conditional subadditivity gives a sum of the slot budgets, not a
single budget. The theorem states its hypothesis in terms of tower errors to
avoid this ambiguity.

\section{E9: Larger observation blocks}
\label{sec:e9}
\begin{enumerate}[label=(\alph*)]
\item A nonempty address $C\subseteq\{1,2,3\}$ meets each of
$\{1,2\},\{1,3\},\{2,3\}$ precisely when $|C|\ge2$.
A singleton misses the complementary pair; every address with at least two
elements meets every pair. Thus the surviving tuple consists of the three
pair slots and the top slot. Applying the two-witness lemma twice, after
observable transfer for each block, gives
\[
 H(U\mid Z_{\{C:|C|\ge2\}})\le3\delta+2\eta_Z.
\]
It does not isolate the top slot alone.
\item For each block $B$, the collection $Z_{\mathcal{J}_B}$ determines
$X_B$: for every $i\in B$ it contains $Z_{\Up\{i\}}$. E6 therefore
extends to $H(U\mid Z_{\mathcal{J}_B})\le H(U\mid X_B)$.
Each $\mathcal{J}_B$ is upward closed. Enumerate
$\mathcal{W}=\{B_1,\ldots,B_N\}$ and intersect the first $j$ contributing
families. Exactly the induction in E8, now with these families as leaves,
gives
\[
 H(U\mid Z_{\mathcal{R}})\le
 \sum_{B\in\mathcal{W}}H(U\mid X_B)+(N-1)\eta_Z,
 \quad \mathcal{R}=\bigcap_{B\in\mathcal{W}}\mathcal{J}_B.
\]
For $N=1$ this is just the transfer inequality. We required a nonempty family
so the induction and the intersection count apply.
\item An address $C$ fails to meet some $(r+1)$-subset of $A$ exactly when
$A\setminus C$ contains at least $r+1$ elements. Thus it meets every such
subset exactly when $|A\setminus C|\le r$. For $r=0$ this says
$A\subseteq C$, recovering the tower. For $r>0$ an address can survive even
if it misses some elements of $A$; the receiver is a larger tuple.
If every witness error is at most $\delta$, write
$N=\binom{|A|}{r+1}$ to obtain the explicit bound
$N\delta+(N-1)\eta_Z$. The restriction $r<|A|$ ensures witnesses exist.
\end{enumerate}

\section{E10: The conditioned score}
\label{sec:e10}
\begin{enumerate}[label=(\alph*)]
\item List $\mathcal{J}_B$ as $A_1,\ldots,A_q$ with strict supersets appearing
before subsets. Such an order exists because strict inclusion has no cycles.
Every strict superset of an active address also meets $B$, so is in this list.
Therefore
\[
 H(Y_{A_t}\mid Y_{\Up A_t\setminus\{A_t\}})
 \ge H(Y_{A_t}\mid Y_{A_1},\ldots,Y_{A_{t-1}}).
\]
Summing and applying the chain rule gives
$\chi(B)\ge H(Y_{\mathcal{J}_B})$, hence $V_B\ge0$.
\item Since $Y_{\mathcal{J}_B}$ determines $X_B$, the joint entropy can be
expanded in two ways:
\[
 H(X_B,Y_{\mathcal{J}_B})=H(Y_{\mathcal{J}_B})
 =H(X_B)+H(Y_{\mathcal{J}_B}\mid X_B).
\]
Subtracting $H(X_B)$ and adding $V_B$ proves
\[
 \chi(B)-H(X_B)=V_B+H(Y_{\mathcal{J}_B}\mid X_B).
\]
\item In E3, the independent model $Y$ has simple score $1+1+1=3$ and
conditioned score $1+1+1=3$. In model $Z$, the simple score is $1+2+2=5$,
while the conditioned score is
\[
 H(C)+H(C,R\mid C)+H(C,T\mid C)=1+1+1=3.
\]
Both models determine $X_1,X_2$, whose joint entropy is three. The conditioned
score charges only the extra information below $C$; its copies add zero
conditional entropy. The simple score charges those copies again.
\item If $A\cap B\ne\varnothing$, every $C\supseteq A$ meets $B$.
Thus $Y_{\Up A}$ is a subtuple of $Y_{\mathcal{J}_B}$ and
\[
 H(Y_{\Up A}\mid X_B)
 \le H(Y_{\mathcal{J}_B}\mid X_B)
 \le\chi(B)-H(X_B),
\]
where the last inequality uses the nonnegative score decomposition.
\end{enumerate}

\subsection{Additional check: equivalence of the perfect definitions}
This argument is available for students who want to connect the two definitions
completely; it is not an extra required exercise.

Let $V_I=\chi(I)-H(Y_{\PI})$. For upward-closed $\mathcal{F},\mathcal{G}$,
partition $\PI$ into
\[
 \mathcal{C}=\mathcal{F}\cap\mathcal{G},\quad
 \mathcal{P}=\mathcal{F}\setminus\mathcal{G},\quad
 \mathcal{Q}=\mathcal{G}\setminus\mathcal{F},\quad
 \mathcal{R}=\PI\setminus(\mathcal{F}\cup\mathcal{G}).
\]
Order these blocks $\mathcal{C},\mathcal{P},\mathcal{Q},\mathcal{R}$, internally
by reverse inclusion. Upward closure ensures this order is valid and no slot
in $\mathcal{Q}$ has a strict superset in $\mathcal{P}$, or conversely.
Conditioning inequalities and chain rules within each block imply
\begin{align*}
 \chi(I)&\ge H(Y_{\mathcal{C}})+H(Y_{\mathcal{P}}\mid Y_{\mathcal{C}})
 +H(Y_{\mathcal{Q}}\mid Y_{\mathcal{C}})
 +H(Y_{\mathcal{R}}\mid Y_{\mathcal{C}},Y_{\mathcal{P}},Y_{\mathcal{Q}})\\
 &=H(Y_{\PI})+I(Y_{\mathcal{P}};Y_{\mathcal{Q}}\mid Y_{\mathcal{C}}).
\end{align*}
The latter mutual information equals
$I(Y_{\mathcal{F}};Y_{\mathcal{G}}\mid Y_{\mathcal{F}\cap\mathcal{G}})$.
Hence $\eta_Y\le V_I$.

If all score gaps vanish, the decomposition forces $V_I=0$, giving the Markov
condition. The singleton gap at $B=\{i\}$ forces the contributing tuple,
and therefore each slot at an address containing $i$, to be determined by
$X_i$. This gives reconstruction.

Conversely, suppose reconstruction and the Markov condition hold.
For every $B$, each active slot $Y_A$ is determined by some $X_i$ with
$i\in A\cap B$, and hence by $X_B$. Therefore
$H(Y_{\mathcal{J}_B}\mid X_B)=0$.
In a reverse-inclusion order, the earlier slots consist of all strict
supersets of the current slot and some incomparable slots. The global
Markov condition implies that conditioning on these extra earlier slots does
not reduce its entropy beyond conditioning on strict supersets: apply the
condition to the upward cone of the current address and the upward-closed
family of earlier addresses. Their intersection is the strict-superset
family. Thus the score's chain-rule inequality is an equality, $V_B=0$.
The decomposition gives $\chi(B)=H(X_B)$ for all $B$.

\section{Further reading}
The correspondence proof specializes Sam Eisenstat's \emph{Condensation: A
Theory of Concepts}, Lemma 6.4 and Theorem 6.8; the exact characterization is
Theorem 5.10. The tower presentation and score decomposition follow the exposition linked in
the student sheet's further reading.
\end{document}
